\documentclass[11pt,a4paper]{article}

\usepackage{amsmath}
\usepackage{amssymb}
\usepackage{amsthm}
\usepackage{booktabs}
\usepackage{array}
\usepackage{longtable}
\usepackage{geometry}
\usepackage{enumitem}
\usepackage{xcolor}
\usepackage[hidelinks]{hyperref}

\geometry{margin=1in}
\emergencystretch=3em
\sloppy

\theoremstyle{plain}
\newtheorem{theorem}{Theorem}
\newtheorem{lemma}[theorem]{Lemma}
\newtheorem{proposition}[theorem]{Proposition}
\theoremstyle{definition}
\newtheorem{definition}[theorem]{Definition}
\theoremstyle{remark}
\newtheorem*{remark}{Remark}

\newcommand{\GD}{G^{D}}
\newcommand{\ED}{E^{D}}
\newcommand{\Cone}{\operatorname{Cone}}
\newcommand{\CC}{\operatorname{CC}}

\title{A Proof Dependency Manifest for a Candidate\\
Computer-Assisted Proof of $f(7,7)=38$}
\author{Jayden Lin\\
\small Northeastern University}
\date{13 September 2026}

\begin{document}
\maketitle

\begin{abstract}
This document is a dependency manifest, not a paper. It states exactly what the
candidate proof of $f(7,7)=38$ consists of, and separates published external
mathematics from project-original arguments, computational reductions, external
database checks, certificate verification, Lean-kernel verification, and the
items still requiring human review. Every claim carries a source and a
verification status. It is written so that a referee can locate the weakest link
quickly.
\end{abstract}

\begin{center}
\fbox{\parbox{0.93\textwidth}{\textbf{Status.} This document describes a
\emph{candidate} computer-assisted proof. All currently known computational
components have been independently checked or formally verified, but the
interpretation of \textsc{akr} Observation~6 / Lemma~7 remains subject to
external human review. The result should not be treated as established.}}
\end{center}

\section*{Disclosure of generative-AI use}
\addcontentsline{toc}{section}{Disclosure of generative-AI use}

Generative-AI systems --- principally Anthropic's Claude, together with OpenAI's
ChatGPT --- were used substantially throughout this project. Their use was not
incidental: it extended to proof exploration and the search for the reduction
strategy; literature-guided reasoning about the published results; generation of
essentially all software, in Python, C++ and Lean; the computational
experimentation itself; debugging; the construction of redundant and independent
verification paths; the Lean formalisation; source auditing of the cited papers;
and the drafting of the project's documentation, including this manifest.

The named author is responsible for the claims made in this document.
AI-generated output was not treated as inherently trustworthy at any point. Each
major claim is discharged by an explicit, separately recorded mechanism ---
source audit against the published papers, independent reimplementation in a
second language, an externally sourced dataset, a machine-checkable certificate
with a standalone checker, or Lean-kernel verification --- and
\S\ref{sec:table} records which mechanism applies to which claim. Several
substantive errors introduced during AI-assisted development were found by those
mechanisms and are catalogued in \S\ref{sec:defects}; they are retained there
rather than quietly repaired, because they calibrate how much checking the
argument has actually received.

Project artifacts and provenance --- source code, generated data, certificates,
build and verification logs, the audit documents, and the complete development
transcripts --- are preserved in the project repository.

No claim is made here about conformance with any particular publisher's or
institution's policy on generative-AI use; that must be assessed separately
against the policy in force.

\section{Claim and current status}

Let $f(m,n)$ denote the maximum number of boundary intersection points of two
simple polygons with $m$ and $n$ sides in general position.

\begin{theorem}[Claimed]\label{thm:main}
$f(7,7)=38$.
\end{theorem}

Three facts frame the problem.

\begin{enumerate}[leftmargin=*]
  \item \textbf{$38$ is realizable.} An explicit pair of integer-coordinate
    heptagons attains $38$ crossings (\S\ref{sec:lean}, \textsf{LB38}); this is
    verified by the Lean kernel.
  \item \textbf{The published state of the art leaves only $40$.}
    {\v C}ern{\'y} et al.~\cite[Theorem 4]{CernyEtAl2003} prove
    $f(k,\ell)\le k\ell-\lceil k/6\rceil-\ell$ for odd $7\le k\le\ell$, which at
    $k=\ell=7$ gives $f(7,7)\le 40$; \cite{AckermanKeszeghRote2022} is
    asymptotically optimal but vacuous at $m=n=7$. Crossings of two closed curves
    in general position are even, so $f(7,7)\in\{38,40\}$ and the whole question
    is whether $40$ is attainable.

    \emph{This bound is context, not a dependency.} Importing it would shorten
    the argument --- $I\le40$, $I\neq40$ and parity would give $I\le38$ at once,
    with no $I=42$ case --- but \cite[Theorem 4]{CernyEtAl2003} rests on
    \cite[Lemma 3]{CernyEtAl2003}, a four-page case analysis this project has not
    audited, and importing it would contradict the project's own accounting of
    its dependencies. The argument here instead uses only
    \cite[Lemma 1]{CernyEtAl2003} (\textsf{C1}), which gives the weaker cap
    $I\le42$, and excludes $I=42$ separately in \S\ref{sec:i42}. {\v C}ern{\'y}
    Theorem 4 is therefore independent corroboration of the upper bound, not an
    input to it.
  \item \textbf{The upper bound is a proof by contradiction.} One assumes a pair
    with $I=40$ and derives a contradiction. Everything in
    \S\ref{sec:flow}--\S\ref{sec:cert} serves that single reductio.
\end{enumerate}

We do \emph{not} claim Theorem~\ref{thm:main} is established. We claim that no
computational gap is currently known, and that the remaining risk is
concentrated in one identified place (\textsf{A7-EXT}, \S\ref{sec:a7ext}).

\section{One-page proof flow}\label{sec:flow}

Write $\GD$ for the bipartite \emph{disjointness graph} on the $7+7$ sides, with
$(p_i,q_j)\in\ED$ iff the closed segments $p_i$ and $q_j$ are disjoint. Since two
non-collinear segments meet at most once, crossings biject with intersecting side
pairs, so
\[
  I + |\ED| = 49 .
\]

The reductio proceeds as follows; the label under each arrow names the
justification.
{\small\[
\begin{array}{rcl}
  I=40
  &\stackrel{\text{\footnotesize $I+|\ED|=49$}}{\Longrightarrow}&
  |\ED| = 9 \\[4pt]
  &\stackrel{\text{\footnotesize \textsf{C1} ({\v C}ern{\'y} L1)}}{\Longrightarrow}&
  \delta(\GD)\ge 1 \\[4pt]
  &\stackrel{\text{\footnotesize \textsf{LA} (Lemma A + parity)}}{\Longrightarrow}&
  P \text{ and } Q \text{ are both \emph{exceptional}} \\[4pt]
  &\stackrel{\text{\footnotesize \textsf{OT},\textsf{EX} (order-type classification)}}{\Longrightarrow}&
  (\text{profile}_P,\text{profile}_Q) \in \{14\}\times\{14\} \\[4pt]
  &\stackrel{\text{\footnotesize \textsf{GR},\textsf{COMP} (enumeration; \textsf{A5-C},\textsf{A7-EXT},\textsf{C2})}}{\Longrightarrow}&
  225 \text{ terminal triples on } 121 \text{ graphs} \\[4pt]
  &\stackrel{\text{\footnotesize \textsf{A5-A},\textsf{CERT},\textsf{LEAN-END}}}{\Longrightarrow}&
  \text{every terminal triple violates \textsc{akr} Lemma 5} \\[4pt]
  &\Longrightarrow& \bot .
\end{array}
\]}

Hence $I\neq 40$; with $I\le42$ (\textsf{C1}), $I\neq42$ (\S\ref{sec:i42}) and
$I$ even, this gives $I\le 38$, and the explicit construction \textsf{LB38}
attains $38$. Therefore $f(7,7)=38$.

\begin{remark}
All of this is conditional on the general-position regime \textsf{GP}, imported
from \cite[\S9.1]{AckermanKeszeghRote2022}; see \S\ref{sec:gp}.
\end{remark}

\section{Dependency table}\label{sec:table}

\newcommand{\stat}[1]{\footnotesize\textsc{#1}}
{\footnotesize\setlength{\tabcolsep}{3pt}
\begin{longtable}{@{}>{\raggedright\arraybackslash}p{1.3cm}>{\raggedright\arraybackslash}p{3.5cm}>{\raggedright\arraybackslash}p{1.8cm}>{\raggedright\arraybackslash}p{2.4cm}>{\raggedright\arraybackslash}p{2.4cm}>{\raggedright\arraybackslash}p{2.8cm}@{}}
\toprule
\textbf{ID} & \textbf{Claim} & \textbf{Type} & \textbf{Source} & \textbf{Verification} & \textbf{Remaining risk} \\
\midrule
\endfirsthead
\toprule
\textbf{ID} & \textbf{Claim} & \textbf{Type} & \textbf{Source} & \textbf{Verification} & \textbf{Remaining risk} \\
\midrule
\endhead
\bottomrule
\endfoot

\textsf{GP} & General-position regime may be assumed w.l.o.g.
  & published theorem & \cite[\S9.1,\S3]{AckermanKeszeghRote2022}
  & \stat{source-audited} & none known; project's own earlier proof was wrong and is removed \\
\addlinespace
\textsf{C1} & A line meets a simple $7$-gon in at most $6$ sides; hence $\delta(\GD)\ge1$
  & published theorem & \cite[Lemma 1]{CernyEtAl2003}
  & \stat{source-audited} & none known \\
\addlinespace
\textsf{C2} & Forbidden local matching on $4$ consecutive sides with $D\in\operatorname{int}\triangle ABC$
  & published theorem & \cite[Lemma 2]{CernyEtAl2003}
  & \stat{source-audited} & none known \\
\addlinespace
\textsf{C2-BRIDGE} & The implemented filter predicate fires iff an \emph{exact induced} perfect matching exists
  & project inference & this project
  & \stat{lean-kernel-checked} & certifies the Lean transcription of the filter, not the Python/C++ source \\
\addlinespace
\textsf{A5-C} & Lemma 5 applied to \emph{consecutive} side pairs in different components
  & published theorem & \cite[Lemma 5]{AckermanKeszeghRote2022}
  & \stat{source-audited} & none known \\
\addlinespace
\textsf{A5-A} & Lemma 5 applied to \emph{avoiding} side pairs
  & published theorem & \cite[Lemma 5]{AckermanKeszeghRote2022}
  & \stat{source-audited}; bridge \stat{lean-kernel-checked} & none known \\
\addlinespace
\textsf{A7-EXT} & Under the intended reading of ``associated''/``is hooking'', AKR Lemma 7 gives: adjacent $H$-vertices have disjoint status sets
  & published theorem $+$ interpretation & \cite[Lemma 7]{AckermanKeszeghRote2022}, \cite{AckermanKeszeghRote2020}
  & \stat{human review pending} & \textbf{principal open item}; Observation 6 unproved and false under the printed definition of ``associated'' \\
\addlinespace
\textsf{A7-BRIDGE} & Non-empty status sets $+$ disjoint status sets across edges $\Rightarrow$ $H$ bipartite
  & project inference & this project
  & \stat{lean-kernel-checked} & none known; conditional on \textsf{A7-EXT} \\
\addlinespace
\textsf{LA} & Lemma A: an uncrossable $4$-set forces $|\ED|\ge10$, parity gives $\ge11$, so $I\le38$
  & project lemma & this project
  & \stat{lean-kernel-checked} (counting core) & geometric premises $h_1,h_2$ remain hypotheses \\
\addlinespace
\textsf{OT} & There are exactly $135$ order types of $7$ points
  & external dataset cross-check & \cite{AichholzerAurenhammerKrasser2002,OrderTypeDatabase}
  & \stat{externally cross-checked} (exact set equality) & none known \\
\addlinespace
\textsf{EX} & $6{,}342$ simple heptagons; $16$ exceptional; $14$ profile classes
  & finite enumeration & this project
  & \stat{independently reimplemented}; $6{,}342$ \stat{externally cross-checked} & none known \\
\addlinespace
\textsf{GR} & Exactly $13{,}003{,}620$ candidate disjointness graphs
  & finite enumeration & this project
  & \stat{independently reimplemented} $+$ closed form & none known \\
\addlinespace
\textsf{COMP} & Filter pipeline down to $121$ graphs / $225$ triples
  & finite enumeration & this project
  & \stat{independently reimplemented} (exact set match) & inherits \textsf{A7-EXT} \\
\addlinespace
\textsf{CERT} & Each of the $225$ triples violates the \textsc{akr}-L5 consequence
  & certificate & this project
  & \stat{certificate-checked} & completeness of the list is not certified \\
\addlinespace
\textsf{LEAN-LB} & The explicit pair has $38$ crossings and is simple
  & Lean formalization & \texttt{lean/}\allowbreak\texttt{LowerBound}\allowbreak\texttt{.lean}
  & \stat{lean-kernel-checked} & simplicity is orientation-level, not Jordan-curve \\
\addlinespace
\textsf{LEAN-LA} & Counting and parity core of Lemma A
  & Lean formalization & \texttt{lean/}\allowbreak\texttt{LemmaA}\allowbreak\texttt{.lean}
  & \stat{lean-kernel-checked} & proved as an implication \\
\addlinespace
\textsf{LEAN-END} & All $225$ cases fail the \textsc{akr}-L5 condition; conclusion $\leftrightarrow$ predicate
  & Lean formalization & \texttt{lean/}\allowbreak\texttt{Endgame}\allowbreak\texttt{.lean}
  & \stat{lean-kernel-checked} & completeness of the $225$ list not formalized \\
\addlinespace
\textsf{LB38} & $f(7,7)\ge38$
  & construction & this project
  & \stat{lean-kernel-checked} & none known \\
\addlinespace
\textsf{UB42} & $I=42$ is impossible
  & project inference & this project, \S\ref{sec:i42}
  & \stat{manually checked} & inherits \textsf{A7-EXT}; \textsf{C1} gives only $I\le42$, so this branch is needed \\
\addlinespace
\textsf{UB40} & $I=40$ is impossible, hence $f(7,7)\le38$
  & composite & \textsf{GP}, \textsf{C1}, \textsf{C2}, \textsf{A5-C}, \textsf{A5-A}, \textsf{A7-EXT}, \textsf{A7-BRIDGE}, \textsf{LA}, \textsf{COMP}, \textsf{UB42}
  & conditional on \textsf{A7-EXT} & \textbf{fails if \textsf{A7-EXT} fails} \\
\end{longtable}
}

\section{External mathematical dependencies}

\subsection{\textsf{GP} --- general position, AKR \S9}\label{sec:gp}

\textbf{Role.} The whole combinatorial argument is carried out in a
general-position regime. \textsf{GP} is what licenses that restriction.

\textbf{Source.} \cite[\S9.1, and \S3]{AckermanKeszeghRote2022}. Their
perturbation removes the relevant degeneracies without decreasing the number of
boundary intersections.

\textbf{Exact consequences used, and nothing stronger.}
\begin{enumerate}[label=(GP\arabic*),leftmargin=*]
  \item both polygons remain simple;
  \item no vertex of one polygon lies on the boundary of the other, so every
    intersection point is interior to two sides --- this is what allows
    intersections to be counted as proper crossings;
  \item no three vertices of $P$ and $Q$ combined are collinear;
  \item no two sides of $P$ and $Q$ combined are parallel;
  \item consequently the supporting line $\ell(q_j)$ of any side of $Q$ contains
    no vertex of $P$ (it passes through two vertices of $Q$, so a vertex of $P$
    on it would violate (GP3)). This is exactly the hypothesis that
    Lemma~A needs, and it is the \emph{only} kind of line at which Lemma~A is
    ever instantiated.
\end{enumerate}

\begin{remark}[A superseded project argument]
Earlier versions of this project proved \textsf{GP} by a project-specific
perturbation rule. That argument is \textbf{incorrect}: its rule for a
vertex-on-vertex degeneracy, ``push the coincident vertex into $\Cone_Q(w)$'',
destroys the intersection point whenever $\Cone_P(v)\subset\Cone_Q(w)$, yielding
$0$ crossings instead of $\ge1$. It also never addressed the case in which one
side participates in several degeneracies at once. It has been removed from the
active proof and replaced by the citation above. See \S\ref{sec:defects}.
\end{remark}

\begin{remark}[Counting model]
\cite[\S9.3]{AckermanKeszeghRote2022} explicitly leaves the variant counting
intersecting \emph{side pairs} rather than intersection \emph{points} as future
research. Everything here is in the intersection-\emph{point} model, and the
result must not be restated in the pairs model.
\end{remark}

\subsection{\textsf{C1} --- {\v C}ern{\'y} Lemma 1}\label{sec:c1}

A line meets a simple $k$-gon an even number of times, hence at most $k-1$ times
for odd $k$ \cite[Lemma 1]{CernyEtAl2003}. Applied to $\ell(p_i)$ against $Q$ and
to $\ell(q_j)$ against $P$, every side misses at least one side of the other
polygon. The single consequence used is
\[
  \delta(\GD)\ \ge\ 1 .
\]
Together with $|\ED|=49-I$ this also gives $I\le42$.

\subsection{Excluding $I=42$}\label{sec:i42}

The manifest previously asserted twice that $I=42$ is ``excluded separately''
without giving the argument. It is this. (This subsection is a
\emph{project-original} argument --- dependency \textsf{UB42} --- placed here
because it consumes \textsf{C1} directly; it is not itself an external result.)

$I=42$ forces $|\ED|=7$. With $\delta(\GD)\ge1$ on $7+7$ vertices and only $7$
edges, every vertex has degree exactly $1$, so $\GD$ is a \emph{perfect
matching}. Its components are therefore the $7$ matched pairs: each component
contains exactly one side of $P$ and one of $Q$. In particular no two sides of
$P$ share a component.

Hence every consecutive pair $e_i=(p_i,p_{i+1})$ is relevant, so
$V(H)=\mathbb{Z}_7$; and for any $i,k$ with $(k-i)\bmod7\notin\{0,1,6\}$ the four
sides $p_i,p_{i+1},p_k,p_{k+1}$ are distinct and hence lie in four distinct
components. So $H$ is exactly the complement of the $7$-cycle: $7$ vertices and
$14$ edges, containing the triangle $e_0e_2e_4$. A triangle is not bipartite,
contradicting \textsf{A7-EXT} together with \textsf{A7-BRIDGE}.

\begin{remark}
The shipped $|\ED|=7$ sweep reproduces this independently:
$5{,}040 \to 14$ (\textsf{A5-C}) $\to 0$ (\textsf{A7-EXT}). The companion
$|\ED|=8$ sweep gives $476{,}280\to1{,}078\to0$, for $I=41$, which parity already
excludes. Both were a by-product of the historical enumeration bug of
Remark~\ref{rem:bug} and are corroborative rather than load-bearing.
\end{remark}

\begin{remark}[Why not route around \textsf{A7-EXT} here]
$I=42$ could instead be excluded by \cite[Lemma 3]{CernyEtAl2003} directly --- the
six matched partners of six consecutive sides of $P$ form precisely the
configuration that lemma forbids --- which would avoid \textsf{A7-EXT} in this
branch. We do not do this. Since the $I\neq40$ branch is already
\textsf{A7-EXT}-gated and the theorem is the conjunction of the two branches, the
conjunction is \textsf{A7-EXT}-gated either way: routing $I\neq42$ around it
reduces the \emph{number of appearances} of the weak link without reducing the
probability that the theorem survives. It would buy presentation at the cost of
importing a further un-audited lemma.
\end{remark}

\subsection{\textsf{C2} --- {\v C}ern{\'y} Lemma 2}

The only consequence used is a local forbidden configuration. Let
$A,B,C,D,E$ be five consecutive vertices of $P$, giving four consecutive sides
$AB,BC,CD,DE$, and suppose $D\in\operatorname{int}\triangle ABC$. Then for no
four sides of $Q$ is the induced subgraph of $\GD$ on those four sides of $P$ and
those four sides of $Q$ \emph{exactly} a perfect matching
\cite[Lemma 2]{CernyEtAl2003}.

The geometric hypothesis depends only on the order type and is tabulated per
exceptional class as the \emph{Lemma-2 windows}. The word \emph{exactly} is
load-bearing: a block containing a perfect matching \emph{plus} an extra edge is
not forbidden, and the implementation was checked not to reject such blocks.

\subsection{\textsf{A5} --- AKR Lemma 5}

\textbf{Statement used.} If two sides $p_a,p_b$ of $P$ that are \emph{consecutive
or avoiding} lie in different components of $\GD$, then there are consecutive
sides $q_i,q_{i\pm1}$ of $Q$ with
\[
  (p_a,q_i),\,(p_b,q_{i\pm1})\in\ED
  \qquad\text{and}\qquad
  (p_a,q_{i\pm1}),\,(p_b,q_i)\notin\ED .
\]

\textbf{\textsf{A5-C}, consecutive use.} Applied to every consecutive pair in
different components; this is the first filter of the pipeline.

\textbf{\textsf{A5-A}, avoiding use.} Two non-adjacent sides are \emph{avoiding}
when neither contains the intersection point of their supporting lines. This is an
order-type predicate, tabulated per exceptional class. \textsf{A5-A} is the final
filter and the one the certificate records.

\textbf{Combinatorial consequence.} Writing $N(\cdot)$ for neighbourhoods in
$\GD$, the displayed conclusion is equivalent to
\begin{equation}\label{eq:adj}
  \exists\, k\in N(p_a)\setminus N(p_b),\quad
  \exists\, k'\in N(p_b)\setminus N(p_a),\qquad
  k' \equiv k\pm 1 \pmod 7 .
\end{equation}
Both directions are immediate by unfolding: $k\in N(p_a)\setminus N(p_b)$ says
exactly $(p_a,q_k)\in\ED$ and $(p_b,q_k)\notin\ED$.

\begin{remark}
This equivalence is \textbf{Lean-kernel-checked}, as
\texttt{akr5Conclusion\_iff\_adjacent} in \texttt{lean/EndgameDefs.lean}, and
depends on no axioms. It matters because it shows the implemented predicate is
neither a weakening nor a strengthening of the published conclusion.
\end{remark}

\begin{remark}
If $\CC(p_a)\neq\CC(p_b)$ then $N(p_a)\cap N(p_b)=\emptyset$, since a shared
neighbour would join them in one component. The set differences in
\eqref{eq:adj} are then the neighbourhoods themselves.
\end{remark}

\subsection{\textsf{A7-EXT} and \textsf{A7-BRIDGE} --- AKR Lemma 7 and bipartiteness of $H$}\label{sec:a7ext}

This dependency was previously recorded as a single item. It bundles two things
of very different character, and is split here:
\[
  \underbrace{\textsf{A7-EXT}}_{\text{published theorem}+\text{interpretation}}
  \;\Longrightarrow\;
  \underbrace{\textsf{A7-BRIDGE}}_{\text{project inference, Lean-checked}}
  \;\Longrightarrow\;
  H \text{ bipartite}.
\]

\textbf{Definition of $H$.} Fix $P$ (symmetrically for $Q$). Vertices are the
relevant consecutive pairs, $V(H)=\{\,i : \CC(p_i)\neq\CC(p_{i+1})\,\}$, and
$i\sim_H k$ iff the two pairs share no side, i.e.\ $(k-i)\bmod 7\notin\{0,1,6\}$,
\emph{and} the four sides $p_i,p_{i+1},p_k,p_{k+1}$ lie in four distinct
components of $\GD$. Condition (b) is imposed directly, so an $H$-edge reproduces
Lemma 7's hypothesis exactly.

For $e\in V(H)$ let $A(e)$ be the set of pairs of consecutive sides of $Q$
associated to $e$ \emph{in the narrow sense} --- carrying the two cross
non-edges, i.e.\ the full conclusion of Lemma 5 --- and put
\[
  S(e)=\{\text{hooking}: \exists\alpha\in A(e),\ I(\alpha)\in\Cone(e)\}
     \cup\{\text{hooked}: \exists\alpha\in A(e),\ I(e)\in\Cone(\alpha)\}.
\]
$S(e)\neq\emptyset$: $e\in V(H)$ is Lemma 5's hypothesis, so $A(e)\neq\emptyset$,
and Lemma 5's dichotomy realises at least one status.

\subsubsection*{\textsf{A7-EXT} --- external, \stat{human review pending}}

\begin{quote}
Under the intended interpretation of ``associated'' and ``is hooking'',
\cite[Lemma 7]{AckermanKeszeghRote2022} yields: if $e\sim_H e'$ then
$S(e)\cap S(e')=\emptyset$.
\end{quote}

This is the one item that remains subject to external human review, for three
reasons, in increasing order of comfort.
\begin{enumerate}[leftmargin=*]
  \item \textsc{akr}'s Observation 6 (the Axis Property) is stated
    \emph{without proof} in both versions, and is \emph{false} under the printed
    definition of ``associated'', which requires only two edges of $\GD$ and not
    the two cross non-edges. An explicit rational counterexample is recorded in
    the project's review note. Observation 6 is used throughout Lemma 7's case
    analysis.
  \item The conference version \cite{AckermanKeszeghRote2020} opens its proof of
    Lemma 7 by \emph{selecting} the associated pairs subject to the cone
    conditions $I(q_{i'},q_{i'\pm1})\in\Cone(p_i,p_{i+1})$, which forces the
    existential reading of ``is hooking'' used above. The journal version deleted
    those clauses, leaving its own next sentence unsupported. This strongly
    indicates the intended reading but does not state it in one place.
  \item \cite[\S10.2]{AckermanKeszeghRote2022}, journal only, reports the
    authors' own independent computer verification of Lemma 7: $158{,}830$
    arrangements of eight lines, $840$ assignments each after symmetry, yielding
    $14$ configurations of which $13$ fail Lemma 5's hooking/hooked pattern. This
    corroborates the \emph{statement}, not the project's interpretation.
\end{enumerate}

\subsubsection*{\textsf{A7-BRIDGE} --- project inference, \stat{lean-kernel-checked}}

\begin{quote}
If every relevant vertex has a non-empty status set, and adjacent vertices have
disjoint status sets, then $H$ is bipartite.
\end{quote}

\emph{Proof.} Let $e\sim_H e'$. Both $S(e),S(e')$ are non-empty and disjoint
subsets of a two-element set, hence distinct singletons. So if
$S(e)=\{\text{hooking},\text{hooked}\}$ then $e$ has no $H$-neighbour at all; every
non-isolated vertex has $|S(e)|=1$; every $H$-edge joins opposite singletons; and
$e\mapsto S(e)$ is a proper $2$-colouring of the non-isolated part, isolated
vertices being irrelevant to $2$-colourability. $\square$

This inference is formalised in \texttt{lean/}\allowbreak\texttt{Lemma7Bridge}\allowbreak\texttt{.lean}
(\S\ref{sec:lean}). The disjointness condition is an explicit \emph{hypothesis};
AKR Lemma 7 is not assumed, not axiomatised, and not named in the file.

\begin{remark}[What the Lean check does and does not buy]
It settles the ``different associated pairs witness different statuses'' worry
completely: a both-status vertex is \emph{proved} isolated, so the scenario in
which Lemma 7 is applied with inconsistent choices on different $H$-edges is
vacuous. It says nothing whatever about whether \textsc{akr}'s published language
licenses \textsf{A7-EXT}. Lean verifies the project's implication; it does not
verify \textsc{akr}'s geometric theorem.
\end{remark}

\section{Project-original mathematical reductions}

\subsection{Lemma A}

\begin{lemma}[Lemma A]\label{lem:A}
Let $S\subseteq\operatorname{sides}(P)$ satisfy
\[
  (\star)\qquad
  \forall \ell,\ \ \ell\cap V(P)=\emptyset
  \ \Longrightarrow\
  \ell \text{ misses at least one member of } S .
\]
Then $|\ED| \ge n + (m-|S|)$.
\end{lemma}

\begin{proof}[Proof sketch]
For each side $q_j$ of $Q$, (GP5) gives $\ell(q_j)\cap V(P)=\emptyset$, so
$(\star)$ applies and yields $p\in S$ with $\ell(q_j)\cap p=\emptyset$, hence
$q_j\cap p=\emptyset$. This is $n$ edges with distinct $Q$-endpoints, all with
$P$-endpoint in $S$. By \textsf{C1} each $p\notin S$ misses some side of $Q$,
giving $m-|S|$ further edges with distinct $P$-endpoints outside $S$. The two
families are disjoint because their $P$-endpoints lie inside versus outside $S$.
\end{proof}

For $m=n=7$ and $|S|=4$ this gives $|\ED|\ge10$; since $I$ is even, $|\ED|=49-I$
is odd, so $|\ED|\ge11$ and $I\le38$. Call $P$ \emph{exceptional} when no
$4$-subset of its sides satisfies $(\star)$. Then
\[
  I=40 \ \Longrightarrow\ P \text{ and } Q \text{ are both exceptional}.
\]

\textbf{Lean boundary.} \texttt{lean/}\allowbreak\texttt{LemmaA}\allowbreak\texttt{.lean} kernel-checks the counting and
the parity upgrade (\texttt{card\_E\_ge}, \texttt{card\_E\_ge\_eleven},
\texttt{crossings\_le\_38}) as an \emph{implication}. The geometric premises
--- that each $q_j$ is disjoint from some member of $S$, and that each
$p\notin S$ is disjoint from some $q_j$ --- remain explicit hypotheses,
discharged in Lean only for the one explicit certificate pair. The
$(\star)$ condition itself is not formalized.

\subsection{Exceptional classification}

All predicates used --- simplicity, crossing sets, the avoiding relation, the
Lemma-2 windows --- are Boolean functions of the order type of the seven
vertices. The classification is therefore finite:
\[
\begin{aligned}
  135 \text{ order types}\ \times\ 360 \text{ cyclic orders}
  &\ \longrightarrow\ 6{,}342 \text{ simple heptagons} \\
  &\ \longrightarrow\ 16 \text{ exceptional}
  \ \longrightarrow\ 14 \text{ classes}.
\end{aligned}
\]
The $6{,}342$ split as $6{,}007$ with an uncrossable $3$-set, $319$ with an
uncrossable $4$-set, and $16$ exceptional. Each of the $14$ classes has
$|M|=5$, exactly five maximal crossing sets all of size $6$, and two avoiding
pairs --- except one class with none.

\begin{remark}[Strict versus loose, and why it matters]\label{rem:strict}
Two readings of ``uncrossable'' must be distinguished.
\[
  \textbf{strict: } (\star) \text{ above}
  \qquad\text{versus}\qquad
  \textbf{loose: } \text{no line at all meets all four closed sides.}
\]
Measured over all $6{,}342$ simple heptagons these give
\[
  \textbf{strict} \rightsquigarrow 16 \text{ exceptional},
  \qquad
  \textbf{loose} \rightsquigarrow 2{,}268 \text{ exceptional},
\]
differing on $2{,}252$ objects. The \emph{strict} reading is the correct one:
Lemma~A is only ever instantiated at lines $\ell(q_j)$, which avoid $V(P)$ by
(GP5). The $2{,}252$ differing objects are exactly those $4$-sets whose only
transversals pass through a vertex of $P$ --- and no $\ell(q_j)$ can be such a
line. Had the loose phrasing been implemented, the finite search would have been
infeasible.
\end{remark}

\section{Computational reduction}\label{sec:comp}

\subsection{The universe of candidate graphs}

There are $7\times 7 = 49$ side pairs, and $I+|\ED|=49$, so $I=40$ forces
$|\ED|=9$. By \textsf{C1} the minimum degree is at least $1$ on both sides. The
number of bipartite graphs on $7+7$ vertices with exactly $9$ edges and minimum
degree at least $1$ is
\[
  \sum_{j=0}^{7}(-1)^j\binom{7}{j}\,[x^9]\bigl((1+x)^{7-j}-1\bigr)^{7}
  \;=\; 13{,}003{,}620 ,
\]
by inclusion--exclusion over uncovered $Q$-sides. This closed form was evaluated
in exact integer arithmetic independently of both enumerations.

\begin{remark}[A historical enumeration bug]\label{rem:bug}
An earlier implementation reported $13{,}484{,}940$ as ``the number of $9$-edge
graphs''. Its recursion never forced the edge budget to be exhausted at the leaf,
so it enumerated $|\ED|\in\{7,8,9\}$:
\[
  5{,}040 \;+\; 476{,}280 \;+\; 13{,}003{,}620 \;=\; 13{,}484{,}940 .
\]
The searched universe was therefore a \emph{superset} of what was needed. No
candidate was missed and the exclusion was not invalidated --- only the reported
figure was wrong, and the $|\ED|=7$ and $|\ED|=8$ cases were swept for free. The
bug nonetheless motivated the independent reimplementation of \S\ref{sec:impl}.
\end{remark}

\subsection{The filter pipeline}

\begin{center}
{\footnotesize
\begin{tabular}{@{}rlll@{}}
\toprule
\textbf{Count} & \textbf{Filter applied next} & \textbf{Justification} & \textbf{Checked by} \\
\midrule
$13{,}003{,}620$ & \textsc{akr} L5, consecutive pairs & \textsf{A5-C} & Python, C++ \\
$42{,}238$       & \textsc{akr} L7 $\Rightarrow$ $H$ bipartite & \textsf{A7-EXT} $+$ \textsf{A7-BRIDGE} & Python, C++ \\
$24{,}990$       & crossing-set compatibility & elementary (order type) & Python, C++ \\
$2{,}300$        & {\v C}ern{\'y} L2, class-dependent & \textsf{C2} & Python, C++ \\
$121$            & \emph{(change of counted object)} & --- & Python, C++ \\
$225$            & \textsc{akr} L5, avoiding pairs & \textsf{A5-A} & Python, C++, certificate, Lean \\
$0$              & --- & --- & --- \\
\bottomrule
\end{tabular}}
\end{center}

\begin{remark}[The $121\to225$ step is not a filter]\label{rem:121}
This is a change in what is being counted, and it should not be read as a
reduction. The figures $13{,}003{,}620$ through $121$ count \emph{graphs}; the
figure $225$ counts \emph{triples} $(\text{class}_P,\text{class}_Q,G)$. The $225$
triples live on exactly those same $121$ graphs: a graph may be compatible with
several exceptional classes on either side, and $225=\sum_G |ok_P(G)|\cdot|ok_Q(G)|$
over the $121$ surviving graphs. Every one of the $121$ graphs has degree sequence
$(2,2,1,1,1,1,1)$ on both sides, and only $9$ of the $14$ classes occur.
\end{remark}

\begin{remark}[{\v C}ern{\'y} Lemma 3 is not a dependency]
An earlier pipeline also applied {\v C}ern{\'y} Lemma 3 between \textsf{A7-EXT} and
the compatibility filter, reducing $24{,}990$ to $9{,}702$. Running the pipeline
with and without it produces \emph{identical} results at every subsequent stage
and an identical set of $225$ terminal triples: the graphs Lemma 3 removes are
all removed by crossing-set compatibility in any case. Lemma 3 is therefore
\emph{not} a logical dependency of this proof and is omitted throughout.
\end{remark}

\section{Independent implementations}\label{sec:impl}

\textbf{Python (reference).} The readable implementation: recursion over
$P$-side neighbourhoods with a running edge budget, union-find components,
DFS $2$-colouring for bipartiteness, a $128$-entry lookup table for
compatibility, and iteration over all $\binom{7}{4}$ subsets for {\v C}ern{\'y}
L2.

\textbf{C++17 (independent reduction).} Written from a written specification,
with every major algorithm chosen differently: enumeration transposed to
$Q$-side columns and partitioned by an explicit degree sequence rather than an
edge budget; graphs keyed as a packed $49$-bit word; components by iterative
bitmask closure; bipartiteness by union-find \emph{with parity}; compatibility by
direct containment in the five maximal crossing sets; {\v C}ern{\'y} L2 by
enumeration of systems of distinct representatives.

\textbf{Agreement.} Both produce the same universe $13{,}003{,}620$; the same
intermediate counts $42{,}238$, $24{,}990$, $2{,}300$, $121$; the same $121$
graphs; and the same canonical set of $225$ terminal triples, with
\emph{zero set difference} in either direction.

\begin{remark}[Scope of the independence claim]
This is \textbf{implementation-independent in algorithm and data representation,
but not fully author- or specification-independent}: the second implementation
was written by an agent that had already seen the first implementation and its
specification. The agreement is therefore strong evidence against transcription
and coding errors, and against a missed branch in either enumeration, but it is
\emph{not} evidence against a shared misreading of the published lemmas.
\end{remark}

\section{Certificates}\label{sec:cert}

A \emph{certificate} here is a file that lets a referee re-derive the terminal
contradictions without trusting --- or rerunning --- the exhaustive search. For
each of the $225$ terminal triples it stores the class identifiers, a canonical
graph encoding (seven neighbourhood bitmasks), the degree sequences, the
component partitions, the relevant side pair and its kind, both neighbourhoods,
the two set differences, and the list of cyclic adjacencies found, which is empty
in every case.

The standalone checker recomputes every quantity from the canonical encoding:
neighbourhoods, degree sequences, components (by breadth-first search, where the
generator used union-find), the set differences, and the adjacency test. The
certificate's convenience fields are treated as claims to be checked, never as
input.

\begin{remark}[Component labels are compared modulo relabelling]
Component identifiers are arbitrary. The checker compares the induced
\emph{partitions}, not the integer labels. An earlier version of the comparison
compared raw labels and reported spurious failures on $217$ of $225$ cases; the
check was corrected to be relabelling-invariant.
\end{remark}

\textbf{Adversarial controls.} The checker was validated against ten deliberate
corruptions, all detected: a case marked as surviving, a pair not avoiding for
its class, falsified neighbourhoods, falsified difference sets, a falsified
adjacency list, a graph edited to break the edge count or minimum degree, a wrong
degree sequence, a wrong component partition, and --- the essential control --- an
\emph{honestly constructed} witness using a different avoiding pair of the same
class for which Lemma 5 genuinely \emph{is} satisfied. That last case is rejected
with a message naming the satisfying adjacency, which confirms the checker does
not simply reject everything.

\subsection*{A worked terminal case}

Take the first certificate entry. The graph is given by the $P$-side
neighbourhoods
\[
  N(p_0),\dots,N(p_6)
  \;=\;
  \{0\},\{0\},\{1\},\{2,4\},\{3\},\{4\},\{5,6\},
\]
with $P$ assigned exceptional class $12$, whose $M$-set is $\{0,2,3,4,6\}$ and
whose avoiding pairs are $(0,2)$ and $(4,6)$. The sides $p_4$ and $p_6$ are
avoiding and lie in different components. Now
\[
  N(p_4)=\{3\},\qquad N(p_6)=\{5,6\},
\]
so $N(p_4)\setminus N(p_6)=\{3\}$ and $N(p_6)\setminus N(p_4)=\{5,6\}$. In
$\mathbb{Z}_7$ the element $3$ is cyclically adjacent only to $2$ and $4$, and
neither lies in $\{5,6\}$. So \eqref{eq:adj} fails, contradicting \textsf{A5-A}.
All $225$ cases have this shape.

For contrast, the \emph{other} avoiding pair $(0,2)$ of the same class on the same
graph gives $N(p_0)\setminus N(p_2)=\{0\}$ and $N(p_2)\setminus N(p_0)=\{1\}$,
and $0$ \emph{is} adjacent to $1$ --- so \eqref{eq:adj} holds there. The
condition is genuinely discriminating: across the $225$ certificate graphs,
$2{,}016$ of the $4{,}725$ side pairs satisfy it.

\section{Lean trust boundary}\label{sec:lean}

Environment: Lean 4.33.1; Mathlib tag \texttt{v4.33.1}, revision
\texttt{0df444a3} \cite{MathlibCommunity}. Build: \texttt{lake build}.

None of the three files contains \texttt{sorry}, \texttt{admit}, a
project-specific \texttt{axiom}, \texttt{unsafe}, or \texttt{native\_decide};
this was verified by scanning the sources with comments stripped. All
\texttt{decide} calls are ordinary kernel reduction.

\subsection*{\texttt{LowerBound.lean} --- \textsf{LEAN-LB}}
\textbf{Imports: none.} Core Lean only; verified to compile with an empty
\texttt{LEAN\_PATH}, so Mathlib is not in its trust base.
\textbf{Axioms:} \texttt{propext} only.

Twenty theorems, including \texttt{P\_simple}, \texttt{Q\_simple} (distinct
vertices; no three of the seven collinear; consecutive sides non-collinear; no
proper crossing of non-adjacent sides), \texttt{general\_position} (no vertex of
one polygon on a side-line of the other), \texttt{crossing\_count\_eq\_38},
\texttt{per\_P\_side\_counts} $=[6,6,6,6,6,6,2]$, \texttt{disjoint\_count\_eq\_11}
and \texttt{crossCount\_add\_disjointCount}. Crucially
\texttt{allPairs\_length} $=49$, \texttt{allPairs\_nodup} and
\texttt{allPairs\_complete} establish that the count of $38$ is taken over
provably all $49$ side pairs.

\textbf{Not formalized:} the step from these orientation-level facts to ``the
closed polygonal curve is a simple closed curve in $\mathbb{R}^2$''.

\subsection*{\texttt{LemmaA.lean} --- \textsf{LEAN-LA}}
\textbf{Imports:} Mathlib (\texttt{Finset.Card}, \texttt{Fintype.Prod},
\texttt{Fintype.Card}, \texttt{Algebra.Ring.Parity}) and \texttt{LowerBound}.
\textbf{Axioms:} \texttt{propext}, \texttt{Classical.choice}, \texttt{Quot.sound}.

Theorems \texttt{card\_E\_ge}, \texttt{card\_E\_ge\_ten},
\texttt{ge\_eleven\_of\_ge\_ten\_of\_odd}, \texttt{card\_E\_ge\_eleven},
\texttt{crossings\_le\_38}, together with a non-vacuity block
(\texttt{Rcert\_h1}, \texttt{Rcert\_h2}, \texttt{Rcert\_card\_E} $=11$,
\texttt{nonvacuous}, \texttt{cert\_I\_le\_38},
\texttt{hypotheses\_are\_restrictive}) showing the hypotheses are simultaneously
satisfiable, the bound tight, and the empty relation excluded.

\begin{remark}[A defect in the earlier skeleton]
An earlier, never-compiled version of this file did not have \texttt{Odd} in
scope. Lean's auto-bound implicits silently turned it into a free variable, so
\texttt{card\_E\_ge\_eleven} and \texttt{crossings\_le\_38} elaborated as vacuous
statements \emph{with no error reported}. The file now sets
\texttt{autoImplicit false}, which makes this class of failure impossible.
\end{remark}

\subsection*{\texttt{Endgame.lean} --- \textsf{LEAN-END}}
\textbf{Imports:} \texttt{EndgameData} only, which imports \texttt{EndgameDefs},
which imports nothing --- core Lean throughout, confirmed standalone with bare
\texttt{lean} and \texttt{LEAN\_PATH=.}\ \textbf{Axioms:} \texttt{propext} and
\texttt{Quot.sound} only; notably \emph{no} \texttt{Classical.choice}.

Principal theorems:
\begin{itemize}[leftmargin=*,nosep]
  \item \texttt{card\_certificates} : the list has length $225$;
  \item \texttt{keys\_distinct} : no duplicate records under the canonical key;
  \item \texttt{all\_225\_cases\_fail\_AKR5\_requirement} and
    \texttt{all\_225\_cases\_fail\_AKR5\_conclusion};
  \item \texttt{all\_225\_pairs\_in\_different\_components} : component
    separation is \emph{recomputed in the kernel} via a separator argument, not
    taken from the stored data;
  \item \texttt{akr5Conclusion\_iff\_adjacent} : the published conclusion and the
    implemented predicate are equivalent; this depends on no axioms at all;
  \item \texttt{akr5\_not\_vacuous} : $2{,}016$ of the $4{,}725$ side pairs in
    these graphs \emph{do} satisfy the condition.
\end{itemize}
The Lean data file is generated deterministically from the committed certificate
JSON, whose hash is embedded in the generated header.

\subsection*{\texttt{Lemma7Bridge.lean} --- \textsf{A7-BRIDGE}}
\textbf{Imports: none.} Core Lean; builds standalone with an empty
\texttt{LEAN\_PATH}. \textbf{Axioms: none at all} --- every theorem reports
\emph{does not depend on any axioms}, not even \texttt{propext}.

Formalises only the implication of \S\ref{sec:a7ext}. The two hypotheses are
\texttt{hne}~:~$\forall e,\exists s,\ S\,e\,s$ and
\texttt{hdisj}~:~$\forall e\,e'\,s,\ \mathrm{Adj}\,e\,e'\to S\,e\,s\to S\,e'\,s\to\bot$.
Theorems: \texttt{adj\_irrefl} (irreflexivity is derived, not assumed),
\texttt{isolated\_of\_both}, \texttt{status\_unique\_of\_adj},
\texttt{statusIs\_of\_adj}, \texttt{opposite\_of\_adj},
\texttt{disjoint\_statuses\_imply\_bipartite} (main), and a partition-form
variant carrying no decidability hypothesis. Non-vacuity is proved rather than
asserted: a concrete model with a real edge; a both-status vertex shown isolated;
and a triangle shown both non-bipartite and admitting no valid status assignment,
so the hypotheses genuinely constrain the graph.

\textbf{AKR Lemma 7 does not appear} --- not as an axiom, not by name. Grepping
the code with comments stripped finds no occurrence of \textsc{akr} or
Observation.

\subsection*{\texttt{CernyL2Interface.lean} --- \textsf{C2-BRIDGE}}
\textbf{Imports: none.} \textbf{Axioms:} the three standard ones; the pure
\texttt{decide} facts are \texttt{propext}-only.

Proves \texttt{l2FiresB\_iff}: the implemented filter fires iff an \emph{exact
induced} perfect matching exists, the single clause
$\mathrm{Edge}\,(s_t)\,(q_u)\leftrightarrow t=u$ carrying both halves --- all four
matching edges present, all twelve cross pairs absent. The decisive fixtures:
a matching plus an extra edge \emph{inside} the $4$-block does \textbf{not} fire,
while a matching plus an edge \emph{outside} the block \textbf{still} fires, so
the predicate is certified to be about the induced subgraph rather than about
degrees.

\begin{remark}[Limits of \textsf{C2-BRIDGE}]
It certifies the Lean transcription of the filter, \emph{not} the Python/C++
source: there is no extraction, and the correspondence remains a human-eyeball
link. One representational shortcut is unverified --- the code computes
$\mathrm{popcount}(N_s\wedge S)$ where the Lean definition sums four indicator
bits; these agree exactly when the four indices are distinct, which the
enumeration enforces, but the bit-level identity is not proved. And it says
nothing about {\v C}ern{\'y} Lemma 2 itself, nor about the per-class Lemma-2
window sets, which remain fully external.
\end{remark}

\begin{center}
\fbox{\parbox{0.93\textwidth}{Key finite and combinatorial components are
formally verified in Lean; the published geometric lemmas themselves remain
external dependencies. In particular $f(7,7)=38$ is \textbf{not} a Lean theorem,
and the completeness of the $225$-case list is not formalized.}}
\end{center}

\section{Trust DAG}

{\footnotesize
\begin{verbatim}
   AKR S9   (GP) ----------+
   Cerny L1 (C1) ----------+
   Cerny L2 (C2) ----------+        Lemma A (LA)  [Lean: counting core]
   AKR L5   (A5) ----------+               |
   AKR L7 (A7-EXT) * ------+               v
                           |     exceptional classification (OT, EX)
                           |       135 order types   [external DB]
                           |       6,342 simple      [external DB]
                           |       16 -> 14 classes  [reimplemented]
                           v               |
                    finite reduction (GR, COMP)
                      /                  \
                 Python                 C++17
                      \                  /
                    225 terminal triples
                   (exact set match, 121 graphs)
                             |
                       certificates (CERT)
                  standalone checker, 10 controls
                             |
                        Lean (LEAN-END)
                 all 225 fail AKR-L5; iff bridge
                             |
                          I != 40
                             |
             explicit 38-example (LB38) [Lean: LEAN-LB]
                             |
                      f(7,7) = 38   (candidate)

   * A7-EXT: HUMAN REVIEW PENDING -- see the A7-EXT subsection of Section 4.  The chain is
         A7-EXT  ==>  A7-BRIDGE  ==>  H bipartite
     where A7-EXT is the external reading of AKR Observation 6 / Lemma 7
     (adjacent H-vertices have disjoint status sets), and A7-BRIDGE is the
     project's own inference from that to bipartiteness -- the latter is
     LEAN-KERNEL-CHECKED, with the disjointness condition as an explicit
     hypothesis and no AKR axiom.
     If A7-EXT fails, the H-bipartite filter is unjustified and the
     reduction from 42,238 onwards must be redone; UB40 fails.
\end{verbatim}
}

\section{Known defects found during development}\label{sec:defects}

Listed because they calibrate how much checking this argument has had.

\begin{enumerate}[leftmargin=*]
  \item \textbf{Enumeration bug ($13{,}484{,}940$).} The leaf test omitted the
    edge-budget condition, so $|\ED|\in\{7,8,9\}$ was enumerated.
    \emph{Effect on the theorem:} none --- a superset was searched, so no
    candidate was missed. \emph{Repair:} corrected leaf test; the buggy file is
    preserved unpatched as evidence. \emph{Superseded by:} the corrected
    enumeration and the closed form of \S\ref{sec:comp}.
  \item \textbf{Raw component-ID comparison.} A verification script compared
    component \emph{labels} rather than the partitions they induce, and reported
    failures on $217$ of $225$ cases. \emph{Effect:} none on the proof; it was a
    checker bug. \emph{Repair:} compare partitions. \emph{Superseded by:} the
    current relabelling-invariant checker.
  \item \textbf{Lean \texttt{Odd} / \texttt{autoImplicit} defect.} In the earlier
    \texttt{LemmaA.lean} skeleton, \texttt{Odd} was out of scope and became a
    free implicit variable, so two theorems elaborated as vacuous statements with
    no error. \emph{Effect:} none, because the file had never been compiled.
    \emph{Repair:} rewritten with \texttt{autoImplicit false}.
    \emph{Superseded by:} the current \texttt{LemmaA.lean}.
  \item \textbf{Step 0 perturbation argument.} The project's own
    general-position proof is incorrect when $\Cone_P(v)\subset\Cone_Q(w)$, where
    its rule destroys the intersection point. \emph{Effect on the theorem:} none,
    because the required statement follows from
    \cite[\S9.1]{AckermanKeszeghRote2022}. \emph{Repair:} removed from the active
    proof. \emph{Superseded by:} \textsf{GP}, \S\ref{sec:gp}.
  \item \textbf{Strict versus loose exceptional condition.} Early phrasing said
    ``no line meets all four sides'', which does not license what the computation
    tests. \emph{Effect if it had been implemented:} severe --- $2{,}268$
    exceptional objects instead of $16$, and an infeasible search.
    \emph{Repair:} the strict, vertex-avoiding form $(\star)$.
    \emph{Superseded by:} Lemma~\ref{lem:A} and Remark~\ref{rem:strict}.
\end{enumerate}

\section{Remaining work}

\begin{enumerate}[leftmargin=*,nosep]
  \item External human review of the \textsc{akr} Observation 6 / Lemma 7
    interpretation (\S\ref{sec:a7ext}).
  \item Ideally, confirmation from one of the \textsc{akr} authors on the
    intended reading of ``associated'' and ``is hooking''.
  \item Independent expert review of the overall proof.
  \item Conversion of this manifest into a readable paper.
  \item The human author learning the argument well enough to defend it.
  \item Further Lean formalization \emph{only} where it improves the trust
    boundary.
\end{enumerate}

Formalizing all of {\v C}ern{\'y} et al.\ or all of \textsc{akr} is
\textbf{not} required by this proof architecture. Those results are used as
published black boxes, and the appropriate check on them is source auditing and
expert review, not formalization.

\section{Final trust statement}

\begin{center}
\fbox{\parbox{0.93\textwidth}{
The project currently contains a heavily audited \emph{candidate}
computer-assisted proof of $f(7,7)=38$. The novel finite reductions and the
terminal elimination have been independently implemented, certificate checked,
and in key places Lean-kernel verified. The principal remaining mathematical
issue is \textsf{A7-EXT}: whether the published Observation 6 / Lemma 7 language
licenses precisely the disjoint-status consequence the argument uses.
Conditional on that consequence, the deduction that the auxiliary graph $H$ is
bipartite (\textsf{A7-BRIDGE}) is formally verified in Lean. Independent expert
review remains required before the result should be treated as established.
}}
\end{center}

\bibliographystyle{plain}
\bibliography{references}

\end{document}
